\documentclass[10pt,twocolumn]{article}

% -----------------------------------------------------------------------
% Packages
% -----------------------------------------------------------------------
\usepackage[margin=1in]{geometry}
\usepackage{amsmath,amssymb,amsthm}
\usepackage{graphicx}
\usepackage{booktabs}
\usepackage{hyperref}
\usepackage{xcolor}
\usepackage{algorithm}
\usepackage{algorithmic}
\usepackage{microtype}
\usepackage{natbib}
\usepackage{caption}

% -----------------------------------------------------------------------
% Custom commands
% -----------------------------------------------------------------------
\newcommand{\tbdsection}[1]{\textcolor{gray}{[TBD: #1]}}
\newcommand{\tbd}[1]{\textcolor{red}{\textbf{[TBD: #1]}}}

% -----------------------------------------------------------------------
% Document metadata
% -----------------------------------------------------------------------
\title{%
  Learning What to Learn:\\
  Adaptive Sample Selection for Efficient Neural Network Training
}

\author{%
  [Author Name]\\
  \textit{Department of Artificial Intelligence, [University]}\\
  \texttt{[email@university.edu]}
}

\date{September 2026}

% -----------------------------------------------------------------------
\begin{document}
% -----------------------------------------------------------------------

\maketitle

% -----------------------------------------------------------------------
\begin{abstract}
% -----------------------------------------------------------------------
Training neural networks on large datasets is increasingly expensive.
Not all training samples contribute equally to learning at every stage
of training. We investigate whether \emph{adaptive sample selection}---
choosing which examples to train on based on per-sample difficulty,
uncertainty, and feature-space diversity---can improve training efficiency
compared to uniform random sampling under a fixed data budget.

We propose a lightweight, forward-pass-only adaptive selector that
combines cross-entropy loss, softmax entropy, and feature-space diversity
with a phase-dependent mixing schedule.
We implement four baselines (random, difficulty-only, uncertainty-only,
diversity-only) and evaluate all methods on Fashion-MNIST~\citep{xiao2017fashionmnist}
using a compact CNN under data budgets of 10\%--100\% of the training pool,
with three random seeds per configuration.

Our experiments suggest \tbd{results pending E001--E005}. We further
present an ablation study isolating the contribution of each selection
signal, and a failure analysis characterising conditions under which
adaptive selection degrades relative to random sampling.

All code, configs, and raw results are publicly available at
\url{https://github.com/[username]/adaptive-sample-selection}.
\end{abstract}

% -----------------------------------------------------------------------
\section{Introduction}
\label{sec:intro}
% -----------------------------------------------------------------------

Modern neural network training can be computationally expensive, even for
moderate-scale datasets.
A standard training run processes every training example once per epoch,
treating all samples as equally informative.
However, empirical evidence suggests this is not the case: some examples
are learned quickly and never forgotten, while others remain persistently
difficult throughout training~\citep{toneva2019forgetting}.

This raises a natural question: \emph{can we select which examples to
include in each training step, and in what order, to make learning
more efficient?}

A long line of work addresses related problems.
Curriculum learning~\citep{bengio2009curriculum} orders examples by
pre-defined difficulty.
Self-paced learning~\citep{kumar2010selfpaced} determines difficulty
dynamically from model loss.
Active learning~\citep{ren2021survey} selects which examples to
\emph{label}; in our setting, labels are available for all data.
Coreset selection~\citep{mirzasoleiman2020craig} finds representative
subsets; our approach is lighter-weight and operates online.

\paragraph{Research question.}
Can lightweight adaptive sample selection improve training efficiency---
measured as validation accuracy per unit of data seen---compared to
uniform random sampling, when all methods operate under the same fixed
data budget?

\paragraph{Contributions.}
We make the following contributions:
\begin{enumerate}
  \item A modular, reproducible framework for comparing sample-selection
    strategies under fixed data budgets.
  \item A controlled empirical study on Fashion-MNIST across five methods,
    five budget levels, and three random seeds.
  \item A lightweight adaptive selector combining difficulty, uncertainty,
    and diversity with a phase-dependent schedule.
  \item An ablation study and failure analysis documenting when adaptive
    selection helps and when it does not.
\end{enumerate}

We do not claim algorithmic novelty over the existing literature.
Our contribution is primarily \emph{empirical}: a rigorous, reproducible
comparison under a CPU-scale setting that is underrepresented in the literature.

% -----------------------------------------------------------------------
\section{Related Work}
\label{sec:related}
% -----------------------------------------------------------------------

\paragraph{Curriculum learning.}
\citet{bengio2009curriculum} showed that presenting examples from easy to
hard improves convergence and generalisation. The curriculum is pre-defined
and not adaptive to the model state.

\paragraph{Self-paced learning.}
\citet{kumar2010selfpaced} introduced SPL, which jointly learns model
parameters and example selection weights using a regularised objective.
\citet{jiang2015self} extended this with a pre-defined curriculum prior
(SPCL). \citet{jiang2014spld} added a diversity constraint (SPLD),
directly motivating our proposed method.

\paragraph{Hard example mining.}
\citet{shrivastava2016ohem} selects highest-loss samples per batch in
detection (OHEM). \citet{lin2017focal} achieves similar effects through
a differentiable loss reweighting (Focal Loss). Both focus on difficulty
without diversity.

\paragraph{Coreset selection.}
\citet{mirzasoleiman2020craig} treats subset selection as submodular
optimisation, approximating the full gradient with a small representative
set. This is powerful but requires gradient computation over the full pool---
expensive for CPU-only settings.

\paragraph{Training dynamics.}
\citet{toneva2019forgetting} showed that many training examples are
redundant: once learned, they are never forgotten and can be removed
without accuracy loss. This motivates our difficulty/loss-based scoring.

\paragraph{Data valuation.}
\citet{ghorbani2019shapley} use Shapley values to assess individual data
point importance, but exact computation is prohibitive. We use tractable
proxies: loss and entropy.

\paragraph{Data pruning.}
\citet{sorscher2022beyond} demonstrate that selecting harder, less-central
examples outperforms random pruning at small data fractions---directly
relevant to our small-budget experiments.

\paragraph{Importance sampling.}
\citet{zhao2015importance} show that non-uniform sampling proportional to
gradient norm yields provably faster convergence, providing theoretical
justification for our loss-based scoring.

Our work differs from active learning~\citep{settles2009active} in that
all labels are available; we select for training efficiency, not annotation
efficiency.

% -----------------------------------------------------------------------
\section{Problem Formulation}
\label{sec:problem}
% -----------------------------------------------------------------------

Let $\mathcal{D} = \{(x_i, y_i)\}_{i=1}^N$ be the full training pool,
and let $f_\theta$ be a neural network with parameters $\theta$.
We define a \emph{data budget} $B \in (0, 1]$ such that each epoch
selects $n = \lfloor B \cdot N \rfloor$ samples from $\mathcal{D}$.

A \emph{selector} $\pi: \mathcal{D} \times f_\theta \times t \to \mathcal{I}$
maps the current model state and training time $t$ to a set of $n$ indices
$\mathcal{I} \subset [N]$, $|\mathcal{I}| = n$.

The training objective is:
\begin{equation}
  \theta^* = \arg\min_\theta \mathbb{E}_{i \sim \pi} [\ell(f_\theta(x_i), y_i)]
\end{equation}
where $\ell$ is cross-entropy loss.

\paragraph{Efficiency metric.}
Our primary measure is validation accuracy $a(s)$ as a function of
cumulative training samples seen $s$. A method $\pi$ is more efficient
than random baseline $\pi_0$ if, for some target accuracy $a^*$:
\begin{equation}
  s_\pi(a^*) < s_{\pi_0}(a^*)
\end{equation}
where $s_\pi(a^*)$ is the number of samples seen to first achieve $a^*$.

% -----------------------------------------------------------------------
\section{Method}
\label{sec:method}
% -----------------------------------------------------------------------

\subsection{Sample Scoring}

For each sample $x_i$ at epoch $t$, we compute three scores using only
forward passes (no gradient computation beyond the training update):

\paragraph{Difficulty $D_i(t)$.}
The normalised per-sample cross-entropy loss:
\begin{equation}
  D_i(t) = \frac{\ell_i(t) - \min_j \ell_j(t)}{\max_j \ell_j(t) - \min_j \ell_j(t) + \epsilon}
\end{equation}
where $\ell_i(t) = -\log p_\theta(y_i | x_i)$ at epoch $t$.
High $D_i$ indicates the model currently finds $x_i$ difficult.

\paragraph{Uncertainty $U_i(t)$.}
The normalised softmax entropy:
\begin{equation}
  U_i(t) = \frac{H_i(t) - \min_j H_j(t)}{\max_j H_j(t) - \min_j H_j(t) + \epsilon}
\end{equation}
where $H_i(t) = -\sum_c p_\theta(c|x_i) \log p_\theta(c|x_i)$.
High $U_i$ indicates the model is uncertain about the class of $x_i$.

\paragraph{Diversity $V_i(t)$.}
The normalised minimum distance from sample $x_i$'s feature representation
to the nearest previously-selected centroid:
\begin{equation}
  V_i(t) = \frac{d_i(t) - \min_j d_j(t)}{\max_j d_j(t) - \min_j d_j(t) + \epsilon}
\end{equation}
where $d_i(t) = \min_{c \in C_{t-1}} \| \phi(x_i) - c \|_2$,
$\phi(x_i)$ is the penultimate-layer feature, and $C_{t-1}$ is the set of
selected sample features from the previous epoch.

\subsection{Adaptive Selection Policy}

The combined score for sample $x_i$ at epoch $t$ is:
\begin{equation}
  S(x_i, t) = \alpha(t) \cdot D_i(t) + \beta(t) \cdot U_i(t) + \gamma(t) \cdot V_i(t)
  \label{eq:score}
\end{equation}
where $\alpha(t), \beta(t), \gamma(t) \geq 0$ are phase-dependent weights
with $\alpha + \beta + \gamma = 1$.

\paragraph{Phase schedule.}
We use a three-phase schedule motivated by curriculum learning:
\begin{align}
  \text{Phase 1 (epochs 1--10):} &\quad (\alpha, \beta, \gamma) = (0.2, 0.2, 0.6) \\
  \text{Phase 2 (epochs 11--20):} &\quad (\alpha, \beta, \gamma) = (0.4, 0.4, 0.2) \\
  \text{Phase 3 (epochs 21+):} &\quad (\alpha, \beta, \gamma) = (0.5, 0.4, 0.1)
\end{align}
Early training emphasises diversity (broad exploration of the feature space).
Late training emphasises difficulty (focus on boundary-clarifying examples).
Ablation experiments (Section~\ref{sec:ablation}) validate this schedule.

\paragraph{Selection rule.}
Select the top-$n$ samples by score $S(x_i, t)$:
\begin{equation}
  \mathcal{I}_t = \operatorname*{arg\,top}_n \{ S(x_i, t) : i \in [N] \}
\end{equation}

\paragraph{Computational cost.}
One forward pass over the pool: $\mathcal{O}(N)$.
Sorting: $\mathcal{O}(N \log N)$.
Diversity scoring (over candidate subset): $\mathcal{O}(n \cdot |\mathcal{C}|)$.
Memory: $\mathcal{O}(N \cdot D)$ for the feature matrix.

Setting score\_every $= 1$ recomputes scores each epoch (highest quality,
highest overhead). Setting score\_every $= k$ reduces overhead by $k\times$
with a slight staleness penalty.

\subsection{Baselines}

We compare against four baselines:
\begin{enumerate}
  \item \textbf{Random} ($\pi_\text{rand}$): uniform sampling without replacement.
  \item \textbf{Difficulty} ($\pi_D$): top-$n$ by $D_i(t)$ alone.
  \item \textbf{Uncertainty} ($\pi_U$): top-$n$ by $U_i(t)$ alone.
  \item \textbf{Diversity} ($\pi_V$): k-means clustering; select nearest-to-centroid from each cluster.
\end{enumerate}

% -----------------------------------------------------------------------
\section{Experimental Setup}
\label{sec:setup}
% -----------------------------------------------------------------------

\subsection{Dataset}
We use Fashion-MNIST~\citep{xiao2017fashionmnist}: 60,000 training images
and 10,000 test images, 28$\times$28 grayscale, 10 balanced classes.
We split training into a pool of 50,000 and a validation set of 10,000
(stratified by class). The test set is used only for final evaluation.

\subsection{Model}
Our primary architecture (SmallCNN) is a compact CNN with two
convolutional blocks (32 and 64 channels) followed by a fully-connected
classifier (~450k parameters). A secondary MLP (~235k parameters) is used
for robustness experiments (Section~\tbd{E010}).

\subsection{Training Protocol}
All methods use identical training protocols:
\begin{itemize}
  \item Optimiser: Adam ($\text{lr}=10^{-3}$, weight decay $10^{-4}$)
  \item Batch size: 128
  \item Epochs: 30
  \item Data augmentation: random horizontal flip (training only)
\end{itemize}

\subsection{Data Budgets}
Budget $B \in \{0.10, 0.25, 0.50, 0.75, 1.00\}$.
At budget $B$, each epoch selects $n = B \times 50\,000$ samples.
All methods receive exactly $n$ samples per epoch.

\subsection{Evaluation}
Three random seeds per configuration (42, 123, 777).
We report mean $\pm$ standard deviation across seeds.
Primary metric: validation accuracy at equal cumulative samples seen.
Secondary metrics: macro F1, wall-clock time, gradient updates.

\subsection{Hardware}
All experiments run on a CPU-only Intel laptop (8 cores, 8.3\,GB RAM).
No GPU is used, ensuring results are reproducible without specialist hardware.

% -----------------------------------------------------------------------
\section{Results}
\label{sec:results}
% -----------------------------------------------------------------------

\tbd{Results pending completion of experiments E001--E005.
Table~\ref{tab:main} and Figures~\ref{fig:curves}--\ref{fig:pareto}
will be populated from results/processed/mean\_std.csv after all runs complete.}

\begin{table}[t]
  \centering
  \caption{Validation accuracy (mean $\pm$ std, 3 seeds) by method and data budget.}
  \label{tab:main}
  \begin{tabular}{lrrrr}
    \toprule
    Method & 10\% & 25\% & 50\% & 100\% \\
    \midrule
    Random      & \tbd{} & \tbd{} & \tbd{} & \tbd{} \\
    Difficulty  & \tbd{} & \tbd{} & \tbd{} & \tbd{} \\
    Uncertainty & \tbd{} & \tbd{} & \tbd{} & \tbd{} \\
    Diversity   & \tbd{} & \tbd{} & \tbd{} & \tbd{} \\
    Adaptive    & \tbd{} & \tbd{} & \tbd{} & \tbd{} \\
    \bottomrule
  \end{tabular}
\end{table}

% -----------------------------------------------------------------------
\section{Ablation Study}
\label{sec:ablation}
% -----------------------------------------------------------------------

\tbd{Ablation results pending E006. Will compare all 7 component combinations
(D only, U only, V only, D+U, D+V, U+V, D+U+V) at budget 25\%.}

% -----------------------------------------------------------------------
\section{Failure Analysis}
\label{sec:failure}
% -----------------------------------------------------------------------

\tbd{Failure analysis pending E007--E009. Will report performance under:
(a) 10\% and 20\% label noise, (b) 5:1 class imbalance, (c) 30\% duplicate data.}

\paragraph{Anticipated failure modes.}
Based on theory and related work, we anticipate the following:

\begin{enumerate}
  \item \textbf{Label noise amplification}: The difficulty selector is expected
    to over-select mislabelled examples (high loss $\not\Rightarrow$ informative).
    The adaptive selector's diversity component may partially mitigate this.
  \item \textbf{Class imbalance}: A purely difficulty-based selector may over-select
    the minority class (harder to classify). Diversity mitigates this by enforcing
    coverage across feature clusters.
  \item \textbf{Redundant data}: The difficulty selector may repeatedly select
    the same hard examples, wasting the budget on redundant samples.
\end{enumerate}

% -----------------------------------------------------------------------
\section{Discussion}
\label{sec:discussion}
% -----------------------------------------------------------------------

\tbd{Discussion pending experimental results.}

\paragraph{If H1 is supported (adaptive $>$ random at small budgets):}
This would confirm that lightweight, forward-pass-only selection provides
a meaningful efficiency advantage over random sampling on Fashion-MNIST.

\paragraph{If H2 is supported (difficulty hurt at 10\% budget):}
This would be a useful negative result, warning practitioners against
naive hard-example mining when the training pool may contain noisy examples
and when diversity is critical at very small budgets.

\paragraph{If H0 is not rejected (no significant advantage):}
This would be a valid scientific result suggesting that for small, balanced
datasets like Fashion-MNIST, adaptive selection overhead does not justify
the gains, and random sampling should be preferred for simplicity.

% -----------------------------------------------------------------------
\section{Limitations}
\label{sec:limitations}
% -----------------------------------------------------------------------

\begin{enumerate}
  \item \textbf{Single dataset}: All primary experiments use Fashion-MNIST.
    Results may not generalise to larger datasets or more complex vision tasks.
  \item \textbf{CPU-only training}: Wall-clock time comparisons are specific
    to our hardware. GPU experiments may show different overhead characteristics.
  \item \textbf{Three seeds}: Statistical power is limited. We report
    mean $\pm$ std but acknowledge that larger-scale replication is desirable.
  \item \textbf{Hyperparameter sensitivity}: The phase schedule weights
    ($\alpha, \beta, \gamma$) are fixed from preliminary ablation. They may
    not be optimal for other datasets or architectures.
  \item \textbf{Tractable proxies}: We use loss and entropy as proxies for
    true data importance. Shapley values or gradient-based methods may provide
    better estimates at higher computational cost.
\end{enumerate}

% -----------------------------------------------------------------------
\section{Conclusion}
\label{sec:conclusion}
% -----------------------------------------------------------------------

\tbd{Conclusion pending experimental results.}

We present a controlled empirical study of lightweight adaptive sample selection
for efficient neural network training under fixed data budgets.
Our proposed method combines per-sample difficulty, uncertainty, and feature-space
diversity signals with a phase-dependent schedule, motivated by the Self-Paced
Learning with Diversity framework~\citep{jiang2014spld} and curriculum learning
principles~\citep{bengio2009curriculum}.

The primary findings of this study will determine whether \tbd{H1 or H0}.

Regardless of outcome, this work establishes a reusable framework and a
reproducible baseline for CPU-scale adaptive sample selection research---
a contribution accessible to undergraduate researchers without specialist hardware.

% -----------------------------------------------------------------------
\bibliographystyle{plainnat}
\bibliography{references}
% -----------------------------------------------------------------------

\end{document}
